\newcommand{\Bin}{\text{Bin}}
\newcommand{\Poi}{\text{Poi}}
\newcommand{\Geo}{\text{Geo}}
\newcommand{\NB}{\text{NB}}
\newcommand{\Hyp}{\text{Hyp}}


\section{Fitting and Estimation}
This section deals with fitting the \sindist{} over various forms of statistical objects -- MSMs (moment summary measures), PDFs, CDFs and ECDFs.

It is to be noted that since the \sindist{} is a location-scale family, investigating its fitting flexibility over standardized/normalized statistical objects is good enough, since it can always adapt to location-scale transformed analogues of the aforesaid objects.

Further, wherever numerical optimization is required, for finiteness of computation, we operate under the following parametric constraints whilst choosing the best candidate \sindist{}, $\Sinu(a,d,s,k)$:
\begin{enumerate}
    \item $a \in [-10^3, 10^3]$
    \item $d \in [10^{-2}, \, 2\times10^3]$
    \item $s \in [10^{-2}, 10^2]$
    \item $k \in [10^{-2}, 10^2]$
\end{enumerate}
an utilize the \texttt{L-BFGS-B} method of the \texttt{optim} function found in the \texttt{R} programming language.

\subsection{MLE, Score equations}
\begin{align*}
    l(a,d,s,k; \utilde x)
    &= \ln \left[ {1 \over d^nA^n(s,k)} \prod_{i=1}^n \sin^k \pi \left( x_i-a \over d \right)^s \right]\\
    &= -n\ln(d)-n\ln(A(s,k)) + k\sum_{i=1}^n \ln \sin \pi\left( x_i-a \over d \right)^s
\end{align*}

Preliminary notation:
\begin{align*}
    A(s,k) &:= \int_0^1\sin^k(\pi z^s) \ \d z\\
    \nabla_s A &= \pi k \int_0^1 \sin^{k-1}(\pi z^s) \cos(\pi z^s) (\ln z) z^s \ \d z\\
    \nabla_k A &= \int_0^1 \sin^k(\pi z^s) \ln \sin(\pi z^s) \ \d z
\end{align*}
The score equations are:
\begin{align*}
    \nabla_a \ell = 0 &\iff \sum_{i=1}^n z_i^{s-1} \cot(\pi z_i^s) = 0 \\
    \nabla_d \ell = 0 & \iff \pi sk \sum_{i=1}^n z_i^s \cot(\pi z_i^s) = -n \\
    \nabla_s \ell = 0 & \iff  \pi k \sum_{i=1}^n {z^s \cos \pi z^s \ln z  \over \sin \pi z^s} = n {\nabla_s A \over A}\\
    \nabla_k \ell = 0 &\iff \sum_{i=1}^n \ln \sin \pi z_i^s = n{\nabla_k A \over A}
\end{align*}
which may not be solved analytically. We can try figuring out computationally their simultaneous solutions, or we may scrap this entire approach and try maximizing the likelihood computationally through other means.

\subsection{MSM on MSM}
We tried fitting the Moment Summary Measures (MSMs) of the standard $\Sinu(s,k)$ family over different target MSMs using a specific but arbitrary choice of the loss function: $L(\vec m , \vec m_{t}) = (\gamma-\gamma_t)^2 + (\kappa - \kappa_t)^2$. Note that since $\Sinu(a,d,s,k)$ is a location-scale family in $(a,d)$, implying that its first two moments can be freely adjusted by modifying $(a,d)$ alone, it was sufficient to test for the skewness-kurtosis adaptability of the distribution -- hence our choice of the standard $\Sinu(s,k)$ family. The results (minimized loss and the choice of parameters) are provided in \autoref{tab:fitpdf}.
\begin{table}[h!]
    \centering
    \begin{tabularx}{\textwidth}{X|X|X}
        \toprule
        Target Density & Hellinger (a,d,s,k) & \texorpdfstring{$L^2$}{L2} (a,d,s,k)\\
        \midrule
        $\N(0,1)$ & 0.00004277289 {\Tiny (-20.542932, 34.985498, 1.304535, 100)} & \\
        $\mathrm{Lognormal}(0,1)$ & 0.000936801 (0.02827762, 1000, 0.08566970, 35.78972) & \\
        $\mathrm{Skew}\N_{\alpha=1} (0,1)$ & 2.189446e-05 (-7.2216216, 24.8516728, 0.5936823, 100) &\\
        \hline
        $\mathrm{Logistic}(0,1)$ & 0.003149183 (-80.275639, 102.222513, 2.882431, 100) & \\
        $\mathrm{Exp}(1)$ & 0.008101119 (0.0005708485, 6.5585684475, 0.0621246472, 1.6766166256) & \\
        $\mathrm{DE}(0,1)$ & 0.01599712 (-8.895071, 16.872288, 1.085768, 15.608238) & \\
        \hline
        $\Gamma(0.1,1)$ & Improper, like below. WHY! & \\
        $\Gamma(100,1)$ & WHY! 1 (1.468279e-07, 1.000001e-01, 1.000000e-01, 1.000000e-01) & \\
        $\Beta(1,1)$ & 0.0030202, (-0.0007780335, 1.0012066335, 0.1004892336, 0.12190192670) & \\
        \hline
        $\mathrm{Cauchy}(0,1)$ & 0.06878475 (-126.211809, 149.303341, 4.141315, 100) & \\
        \hline
        $\chi^2(1)$ & BAD LIKE GAMMA! & \\
        $\chi^2(10)$ & 5.791378e-06 (-0.3201950, 97.1337699, 0.2815837, 22.0271292) & \\
        $\mathrm{t}_1(0,1)$ & 0.06878477 (-126.114092, 149.207078, 4.138119, 100) & \\
        $\mathrm{t}_{10}(0,1)$ & INTEGRAL DIVERGENT & \\
        $\mathrm{F}(1,1)$ & BAD OUTPUT & \\
        \bottomrule
    \end{tabularx}
    \caption{Some common target densities and the best candidate \sindist{}}
    \label{tab:fitpdf}
\end{table}



\subsection{PDF on PDF}
We tried fitting the PDF of $\Sinu(a,d,s,k)$ over different target PDFs (standard location-scale members only) using numerical optimization and 2 different choices of the loss function (Hellinger distance and $L^2$ distance). The results (minimized distance with the choice of parameters) are provided in \autoref{tab:fitpdf}.
\begin{table}[h!]
    \centering
    \begin{tabularx}{\textwidth}{X|X|X}
        \toprule
        Target Density & Hellinger (a,d,s,k) & \texorpdfstring{$L^2$}{L2} (a,d,s,k)\\
        \midrule
        $\N(0,1)$ & 0.00004277289 {\Tiny (-20.542932, 34.985498, 1.304535, 100)} & \\
        $\mathrm{Lognormal}(0,1)$ & 0.000936801 (0.02827762, 1000, 0.08566970, 35.78972) & \\
        $\mathrm{Skew}\N_{\alpha=1} (0,1)$ & 2.189446e-05 (-7.2216216, 24.8516728, 0.5936823, 100) &\\
        \hline
        $\mathrm{Logistic}(0,1)$ & 0.003149183 (-80.275639, 102.222513, 2.882431, 100) & \\
        $\mathrm{Exp}(1)$ & 0.008101119 (0.0005708485, 6.5585684475, 0.0621246472, 1.6766166256) & \\
        $\mathrm{DE}(0,1)$ & 0.01599712 (-8.895071, 16.872288, 1.085768, 15.608238) & \\
        \hline
        $\Gamma(0.1,1)$ & Improper, like below. WHY! & \\
        $\Gamma(100,1)$ & WHY! 1 (1.468279e-07, 1.000001e-01, 1.000000e-01, 1.000000e-01) & \\
        $\Beta(1,1)$ & 0.0030202, (-0.0007780335, 1.0012066335, 0.1004892336, 0.12190192670) & \\
        \hline
        $\mathrm{Cauchy}(0,1)$ & 0.06878475 (-126.211809, 149.303341, 4.141315, 100) & \\
        \hline
        $\chi^2(1)$ & BAD LIKE GAMMA! & \\
        $\chi^2(10)$ & 5.791378e-06 (-0.3201950, 97.1337699, 0.2815837, 22.0271292) & \\
        $\mathrm{t}_1(0,1)$ & 0.06878477 (-126.114092, 149.207078, 4.138119, 100) & \\
        $\mathrm{t}_{10}(0,1)$ & INTEGRAL DIVERGENT & \\
        $\mathrm{F}(1,1)$ & BAD OUTPUT & \\
        \bottomrule
    \end{tabularx}
    \caption{Some common target densities and the best candidate \sindist{}}
    \label{tab:fitpdf}
\end{table}

\subsection{CDF on CDF}
We tried fitting the CDF of $\Sinu(a,d,s,k)$ over different target CDFs (standard location-scale members only) using numerical optimization and 2 different choices of the loss function. The results (minimized distance with the choice of parameters) are provided in \autoref{tab:fitcdf}.
\begin{table}[h!]
    \centering
    \begin{tabularx}{\textwidth}{X|X|X}
        \toprule
        Target CDF & Hellinger (a,d,s,k) & \texorpdfstring{$L^2$}{L2} (a,d,s,k)\\
        \midrule
        $\Bin(10,0.1)$ & & \\
        $\Bin(10,0.9)$ & & \\
        $\Poi(1)$ & & \\
        $\Poi(5)$ & & \\
        $\Geo(0.1)$ & & \\
        $\Geo(0.9)$ & & \\
        $\NB(10,0.1)$ & & \\
        $\NB(10,0.9)$ & & \\
        $\Hyp(20,16,3)$ & & \\
        $\Hyp(20,6,4)$ & & \\
        \bottomrule
    \end{tabularx}
    \caption{Some common target CDFs and the best candidate \sindist{}}
    \label{tab:fitcdf}
\end{table}

\subsection{Quantile-Density Based Metrics}
Following \cite{staudte2017dqf}, we construct the notion of Divergences based on the density-quantile function $dq(p) = f(F^{-1}(p)) \quad \forall p \in [0,1]$. The advantages of calculating divergence metrics on this transformed domain is that divergences become location-scale invariant and invariant to ``flat'' regions in the CDF.

\section{Divergence Density/Potential}
For some target density $g$, if $D_g(f) = \int_{x \in \mathcal X} L(f(x),g(x)) \d x$ denotes the divergence of candidate density $f$ from target $g$, then define:
\[V(x,y) = L(y, g(x))\]
as the divergence density/potential associated with $g$ and induced by $L$. It gives a glimpse at the topography of penalization of $D(\cdot, \cdot)$ with respect to $g$. This is a useful metric if we want to compare how different tuning parameter $\alpha$'s affect the DPD.
\begin{figure}[h!]
    \centering
    \includegraphics[width=0.7\linewidth]{fig/3-divpot.png}
    \caption{Divergence Potential associated with $\Sinu(1,1)$ and induced by the Hellinger distance}
    \label{fig:divpot}
\end{figure}
The process of fitting a candidate distribution over a given target then becomes merely an act of minimizing the line integral of potential function along the density line whilst ensuring that the area under the density line is 1.

We can conclude that the \sindist{} is indeed more flexible if for different target densities, which may even come from data, it can minimize the divergence better than conventional distributions. At the same time, excess efficiency implies lack of robustness in the presence of contamination. In such cases, we justify the utility of the \sindist{} by saying that it fits to data more flexibly than other distributions, and transparently displays the nature of the underlying divergence metric without interference from the rigidity of its shape.

An apparent benefit of the Divergence Potential over the standard choice of Influence Function in Robust Statistics literature is that it provides more information about the divergence functional than possible through incorporation of a degenerate contamination and subsequent analysis of the limit for an infinitesimal perturbation of that contamination.
